\documentclass[11pt]{article}
\usepackage[review]{acl}

\usepackage{times}
\usepackage{latexsym}
\usepackage[T1]{fontenc}
\usepackage[utf8]{inputenc}
\usepackage{microtype}
\usepackage{inconsolata}
\usepackage{graphicx}
\usepackage{booktabs}
\usepackage{amsmath}
\usepackage{amsfonts}
\usepackage{url}

\title{Adaptive Halting and Evolutionary Latent Control in Recursive Sequence Models}

\author{Anonymous Authors}

\begin{document}
\maketitle

\begin{abstract}
% TODO: Write after finalizing the main text.
% State the task/model setting, causal halting result, EC mechanism finding,
% confirmation boundary, and key limitations without overclaiming.
\end{abstract}

\section{Introduction}
% TODO: Motivation, research gap, RQ1--RQ3, contributions.

\section{Related Work}
% TODO: Adaptive computation; recursive/iterative reasoning;
% evolutionary optimization for neural control; evaluation discipline.

\section{Method}
\subsection{Recursive GRU Seq2Seq Model}
\subsection{Causal Token-level Adaptive Halting}
\subsection{Actual Computation Accounting}
\subsection{Evolutionary Latent Step-size Control}
\subsection{Joint Control}
\subsection{Statistical Criteria}

\section{Experiments}
\subsection{Experimental Setup}
\subsection{Causal Execution and Implementation Validation}
\subsection{Adaptive Halting Results}
\subsection{Fixed-depth Evolutionary Latent Control}
\subsection{Joint Control and Confirmation}
\subsection{Ablations and Mechanistic Analysis}
\subsection{Negative Results and Deviations}

\section{Conclusion}
% TODO: Answer RQ1--RQ3; methodological implications; future work.

% Add a separate Limitations section here if required by the target venue.
% \section{Limitations}

\bibliography{references}
\bibliographystyle{acl_natbib}

\appendix
\section{Additional Experimental Details}
% TODO: Full audit checks, deviations, extra tables, reproducibility details.

\end{document}
